\section{An explicit first-run binary code}\label{sec:benchmark}
\begin{proposition}[Preparation-dependent Gaussian calibration]\label{prop:binary}
Take hidden states $\pm1$, reports $Y\mid i\sim N(i,1)$, transition matrix with diagonal $3/4$ and off-diagonal $1/4$, and initial probability $u_0=1/2$. The scalar predictive coordinate is the probability of the positive hidden state. Let $\Phi$ denote the standard Gaussian distribution function and put
\begin{equation}\label{eq:binary-G}
 G(q)=\frac14+\frac12\{q\Phi(1+\ell)+(1-q)\Phi(1-\ell)\},
 \qquad \ell=\tfrac12\log\frac q{1-q}.
\end{equation}
For each $s\in(0,1]$, the equation
\begin{equation}\label{eq:binary-root}
 q_s=sG(1/2)+(1-s)G(q_s)
\end{equation}
has a unique solution on $[1/2,3/4]$. A homogeneous externally stopped three-label program attains the exact categorical Brier risk
\begin{equation}\label{eq:binary-risk}
 R_s=2q_s(1-q_s).
\end{equation}
Five labels implement the same law with internal halting. The code and its risk are explicit; no claim is made that three or five labels are globally optimal for this benchmark.
\end{proposition}
\begin{proof}
Use the initial label with programmed belief $1/2$, and two recurrent labels with beliefs $q$ and $1-q$. At an incoming belief $p$ send the report to the upper label precisely when $y\ge-\tfrac12\log(p/(1-p))$. This is the common cut at posterior probability $1/2$ after the transition. Symmetry gives equal mass $1/2$ to the two labels in the terminal ensemble. For a symmetric incoming ensemble with means $q,1-q$, direct Gaussian integration gives conditional upper output mean $G(q)$; the first report from preparation $1/2$ gives $G(1/2)$. The resolvent calibration is therefore exactly \eqref{eq:binary-root}, not the invariant equation $q=G(q)$.

Write $\phi$ for the Gaussian density. The identity $q\phi(1+\ell)=(1-q)\phi(1-\ell)$ cancels the derivative of the threshold and yields
\[
 G'(q)=\tfrac12[\Phi(1+\ell)-\Phi(1-\ell)],\qquad
 0\le G'(q)\le\tfrac12.
\]
The map in \eqref{eq:binary-root} preserves $[1/2,3/4]$ and contracts it by at most $(1-s)/2$. Its unique fixed point identifies the actual stopped conditional means. The categorical loss is $2p(1-p)$, which proves \eqref{eq:binary-risk}. Iteration $q_{k+1}=sG(1/2)+(1-s)G(q_k)$ has the rigorous a posteriori error bound
\[
 |q_k-q_s|\le\frac{[(1-s)/2]^k}{1-(1-s)/2}|q_1-q_0|.
\]
The internal construction duplicates the two output labels into continuing and halting copies and counts its initial label.

For explicit geometry constants use $D=[1/8,7/8]$, $K=[1/4,3/4]$, and $A=2+\log7$. The log-odds density is a mixture of variance-four Gaussians with means of magnitude at most $A$. Since the derivative of inverse logistic followed by the transition is at most $8e^{|z|}$, a scalar acquired-density upper is
\[
 \Lambda=\frac{4}{\sqrt{2\pi}}e^{A+2}.
\]
On $|z|\le\log3$ the mixture density is at least $(2\sqrt{2\pi})^{-1}\exp[-(A+\log3)^2/8]$, and the inverse Jacobian is at least $8$. Consequently every row dominates $\epsilon$ times uniform probability on $[3/8,5/8]$, where
\[
 \epsilon=\frac1{\sqrt{2\pi}}\exp[-(A+\log3)^2/8]>0.
\]
These are conservative explicit raw constants, valid without truncation. Scalar versus probability-hyperplane norms change fixed factors only. Thus the benchmark verifies the center theorem as well as giving a concrete table. A full-history checkpoint with two terminal labels can simulate this table's two-valued terminal readout, so its optimal risk is no larger than $R_s$; no unsupported formula for its optimal two-center partition is asserted.
\end{proof}
